% GENERATED FILE. Edit canonical lessons, then run npm run generate:books.
% canonical-source-hash: fnv1a64:dabbe64e311e92cc
% canonical-lessons: RU-C329-kipyatit, RU-C329-kipet, RU-C329-otvarivat, RU-C329-kryshka, RU-C329-nakryvat

\chapter{Boiling}
\label{ch:ru-boiling}
\hlchaptermodality{329}

\begin{chapteropening}
\textbf{By the end of this chapter:} \emph{I can describe boiling: bringing water to the boil, water boiling, cooking food in water, a lid and covering a pan.}

Complete the last lesson of chapter 329: I can describe boiling: bringing water to the boil, water boiling, cooking food in water, a lid and covering a pan.
\end{chapteropening}

\section[\texorpdfstring{\ru{кипятить} (kipyatít)}{кипятить (kipyatít)}]{\ru{кипятить} (kipyatít) — to bring to the boil, to boil}

\label{lesson:RU-C329-kipyatit}



\begin{quote}
Before the new one: say the Russian for to grate, then the Russian for a grater.
\end{quote}

\subsection*{You'll want to know: \ru{кипятить}}
\textbf{\ru{кипятить}} (\emph{kipyatít}) — \textquotedblleft{}to bring to the boil, to boil\textquotedblright{}.

\begin{quote}
\ru{Бабушка} \ru{кипятит} \ru{молоко}.
\end{quote}

\begin{quote}
\emph{Bábushka kipyatít molokó.}
\end{quote}

\begin{quote}
Grandma is boiling milk.
\end{quote}

Imperfective; perfective \emph{\ru{вскипятить}} (\emph{vskipyatít}), \emph{\ru{кипячу}}. \emph{\ru{Кипяток}} — boiling water.

\begin{cousinweb}
\emph{\ru{кипятить}} is built on \emph{\ru{кипеть}} \textquotedblleft{}to boil\textquotedblright{}, inherited from Proto-Slavic \emph{kypěti} (EtymDB): to make something boil.
\end{cousinweb}

\begin{grammarlens}[title={What you've built}]
One more word for telling, explaining and describing at B1.
\end{grammarlens}

\subsection*{Your turn}
\begin{itemize}
\raggedright
  \item \emph{Say it:} \emph{kipyatít}
  \item \emph{Say it:} \emph{kipyatít}, once more
  \item \emph{Say it:} a sentence of your own with \emph{kipyatít}
  \item \emph{Recall:} say the Russian for to grate, then the Russian for a grater, then say \emph{kipyatít} again
\end{itemize}

\begin{tcolorbox}[breakable,colback=teal!4,colframe=teal!35!black,title={Before you move on}]
What does \ru{кипятить} mean? (\textquotedblleft{}To bring to the boil, to boil\textquotedblright{}.) Say it once more.
\end{tcolorbox}
\section[\texorpdfstring{\ru{кипеть} (kipét)}{кипеть (kipét)}]{\ru{кипеть} (kipét) — to boil (of a liquid)}

\label{lesson:RU-C329-kipet}



\begin{quote}
Before the new one: say the Russian for a grater, then the Russian for to bring to the boil.
\end{quote}

\subsection*{You'll want to know: \ru{кипеть}}
\textbf{\ru{кипеть}} (\emph{kipét}) — \textquotedblleft{}to boil (of a liquid)\textquotedblright{}.

\begin{quote}
\ru{Вода} \ru{уже} \ru{кипит}.
\end{quote}

\begin{quote}
\emph{Vodá uzhé kipít.}
\end{quote}

\begin{quote}
The water is already boiling.
\end{quote}

Imperfective; perfective \emph{\ru{вскипеть}} (\emph{vskipét}). \emph{\ru{Кипятить}} boils something; \emph{\ru{кипеть}} is the liquid itself boiling.

\begin{cousinweb}
\emph{\ru{кипеть}} is inherited from Proto-Slavic \emph{kypěti} \textquotedblleft{}to boil\textquotedblright{}, from PIE \emph{kwep-} \textquotedblleft{}to smoke, to boil, to move violently\textquotedblright{} (EtymDB).
\end{cousinweb}

\begin{grammarlens}[title={What you've built}]
One more word for telling, explaining and describing at B1.
\end{grammarlens}

\subsection*{Your turn}
\begin{itemize}
\raggedright
  \item \emph{Say it:} \emph{kipét}
  \item \emph{Say it:} \emph{kipét}, once more
  \item \emph{Say it:} a sentence of your own with \emph{kipét}
  \item \emph{Recall:} say the Russian for a grater, then the Russian for to bring to the boil, then say \emph{kipét} again
\end{itemize}

\begin{tcolorbox}[breakable,colback=teal!4,colframe=teal!35!black,title={Before you move on}]
What does \ru{кипеть} mean? (\textquotedblleft{}To boil (of a liquid)\textquotedblright{}.) Say it once more.
\end{tcolorbox}
\section[\texorpdfstring{\ru{отваривать} (otvárivat)}{отваривать (otvárivat)}]{\ru{отваривать} (otvárivat) — to cook in water, to boil (food)}

\label{lesson:RU-C329-otvarivat}



\begin{quote}
Before the new one: say the Russian for to bring to the boil, then the Russian for to boil.
\end{quote}

\subsection*{You'll want to know: \ru{отваривать}}
\textbf{\ru{отваривать}} (\emph{otvárivat}) — \textquotedblleft{}to cook in water, to boil (food)\textquotedblright{}.

\begin{quote}
\ru{Я} \ru{отварил} \ru{картошку}.
\end{quote}

\begin{quote}
\emph{Ya otvaríl kartóshku.}
\end{quote}

\begin{quote}
I boiled the potatoes.
\end{quote}

Imperfective; perfective \emph{\ru{отварить}} (\emph{otvarít}). \emph{\ru{Варить}} — to cook, to boil in general.

\begin{cousinweb}
\emph{\ru{отваривать}} is built on \emph{\ru{от}-} and \emph{\ru{варить}} \textquotedblleft{}to boil\textquotedblright{}, inherited from Proto-Slavic \emph{variti}, which EtymDB traces to PIE \emph{wer-} \textquotedblleft{}to burn\textquotedblright{}.
\end{cousinweb}

\begin{grammarlens}[title={What you've built}]
One more word for telling, explaining and describing at B1.
\end{grammarlens}

\subsection*{Your turn}
\begin{itemize}
\raggedright
  \item \emph{Say it:} \emph{otvárivat}
  \item \emph{Say it:} \emph{otvárivat}, once more
  \item \emph{Say it:} a sentence of your own with \emph{otvárivat}
  \item \emph{Recall:} say the Russian for to bring to the boil, then the Russian for to boil, then say \emph{otvárivat} again
\end{itemize}

\begin{tcolorbox}[breakable,colback=teal!4,colframe=teal!35!black,title={Before you move on}]
What does \ru{отваривать} mean? (\textquotedblleft{}To cook in water, to boil (food)\textquotedblright{}.) Say it once more.
\end{tcolorbox}
\section[\texorpdfstring{\ru{крышка} (krýshka)}{крышка (krýshka)}]{\ru{крышка} (krýshka) — a lid}

\label{lesson:RU-C329-kryshka}



\begin{quote}
Before the new one: say the Russian for to boil, then the Russian for to cook in water.
\end{quote}

\subsection*{You'll want to know: \ru{крышка}}
\textbf{\ru{крышка}} (\emph{krýshka}) — \textquotedblleft{}a lid\textquotedblright{}.

\begin{quote}
\ru{Крышка} \ru{от} \ru{кастрюли} \ru{упала} \ru{на} \ru{пол}.
\end{quote}

\begin{quote}
\emph{Krýshka ot kastryúli upála na pol.}
\end{quote}

\begin{quote}
The saucepan lid fell on the floor.
\end{quote}

Feminine; genitive plural \emph{\ru{крышек}}.

\begin{cousinweb}
\emph{\ru{крышка}} is built on \emph{\ru{крыть}} \textquotedblleft{}to cover\textquotedblright{}, an inherited Slavic verb, with the ending \emph{-\ru{шка}}: a little cover. EtymDB lists no parent.
\end{cousinweb}

\begin{grammarlens}[title={What you've built}]
One more word for telling, explaining and describing at B1.
\end{grammarlens}

\subsection*{Your turn}
\begin{itemize}
\raggedright
  \item \emph{Say it:} \emph{krýshka}
  \item \emph{Say it:} \emph{krýshka}, once more
  \item \emph{Say it:} a sentence of your own with \emph{krýshka}
  \item \emph{Recall:} say the Russian for to boil, then the Russian for to cook in water, then say \emph{krýshka} again
\end{itemize}

\begin{tcolorbox}[breakable,colback=teal!4,colframe=teal!35!black,title={Before you move on}]
What does \ru{крышка} mean? (\textquotedblleft{}A lid\textquotedblright{}.) Say it once more.
\end{tcolorbox}
\section[\texorpdfstring{\ru{накрывать} (nakryvát)}{накрывать (nakryvát)}]{\ru{накрывать} (nakryvát) — to cover; to lay (a table)}

\label{lesson:RU-C329-nakryvat}



\begin{quote}
Before the new one: say the Russian for to cook in water, then the Russian for a lid.
\end{quote}

\subsection*{You'll want to know: \ru{накрывать}}
\textbf{\ru{накрывать}} (\emph{nakryvát}) — \textquotedblleft{}to cover; to lay (a table)\textquotedblright{}.

\begin{quote}
\ru{Он} \ru{накрыл} \ru{сковороду} \ru{крышкой}.
\end{quote}

\begin{quote}
\emph{On nakrýl skovorodú krýshkoy.}
\end{quote}

\begin{quote}
He covered the pan with a lid.
\end{quote}

Imperfective; perfective \emph{\ru{накрыть}} (\emph{nakrýt}), \emph{\ru{накрою}}. \emph{\ru{Накрыть} \ru{на} \ru{стол}} — to lay a table for a meal.

\begin{cousinweb}
\emph{\ru{накрывать}} is built on \emph{\ru{на}-} \textquotedblleft{}on\textquotedblright{} and \emph{\ru{крыть}} \textquotedblleft{}to cover\textquotedblright{}, an inherited Slavic verb; EtymDB lists no parent for the compound.
\end{cousinweb}

\begin{grammarlens}[title={What you've built}]
One more word for telling, explaining and describing at B1.
\end{grammarlens}

\subsection*{Your turn}
\begin{itemize}
\raggedright
  \item \emph{Say it:} \emph{nakryvát}
  \item \emph{Say it:} \emph{nakryvát}, once more
  \item \emph{Say it:} a sentence of your own with \emph{nakryvát}
  \item \emph{Recall:} say the Russian for to cook in water, then the Russian for a lid, then say \emph{nakryvát} again
\end{itemize}

\begin{tcolorbox}[breakable,colback=teal!4,colframe=teal!35!black,title={Before you move on}]
What does \ru{накрывать} mean? (\textquotedblleft{}To cover; to lay (a table)\textquotedblright{}.) Say it once more.
\end{tcolorbox}
